 \documentclass[a4paper,12pt]{article}
 %%%%%%%%%%%%%%%%%%%%%%%%%%%%%%%%%%%%%%%%%%%%%%%%%%%%%%%%%%%%%%%%%%%%%%%%%%%%%%%%%%%%%%%%%%%%%%%%%%%%%%%%%%%%%%%%%%%%%%%%%%%%%%%%%%%%%%%%%%%%%%%%%%%%%%%%%%%%%%%%%%%%%%%%%%%%%%%%%%%%%%%%%%%%%%%%%%%%%%%%%%%%%%%%%%%%%%%%%%%%%%%%%%%%%%%%%%%%%%%%%%%%%%%%%%%%
 \usepackage{eurosym}
 \usepackage{vmargin}
 \usepackage{amsmath}
 \usepackage{multicol}
 \usepackage{graphics}
 \usepackage{enumerate}
 \usepackage{epsfig}
 \usepackage{framed}
 \usepackage{subfigure}
 \usepackage{fancyhdr}
 
 \setcounter{MaxMatrixCols}{10}
 %TCIDATA{OutputFilter=LATEX.DLL}
 %TCIDATA{Version=5.00.0.2570}
 %TCIDATA{<META NAME="SaveForMode" CONTENT="1">}
 %TCIDATA{LastRevised=Wednesday, February 23, 2011 13:24:34}
 %TCIDATA{<META NAME="GraphicsSave" CONTENT="32">}
 %TCIDATA{Language=American English}
 
 %\pagestyle{fancy}
 \setmarginsrb{20mm}{0mm}{20mm}{25mm}{12mm}{11mm}{0mm}{11mm}
 %\lhead{MA4413 2013} \rhead{Mr. Kevin O'Brien}
 %\chead{Midterm Assessment 1 }
 %\input{tcilatex}
 
 \begin{document}
 	
 	\begin{center}
 		\includegraphics[scale=0.60]{shieldtransparent2}
 	\end{center}
 	
 	\begin{center}
 \vspace{1cm}
 \large \bf {FACULTY OF SCIENCE AND ENGINEERING} \\[0.5cm]
 \normalsize DEPARTMENT OF MATHEMATICS AND STATISTICS \\[1.25cm]
 \large \bf {END OF SEMESTER EXAMINATION} \\[1.5cm]
 	\end{center}
 	
 	\begin{tabular}{ll}
 		MODULE CODE: MA4702 & SEMESTER: Repeat 2016\\[1cm]
 		MODULE TITLE: Technology Mathematics 2 & DURATION OF EXAM: 2.5 hours \\[1cm]
 		LECTURER: Kevin O'Brien & GRADING SCHEME: 100 marks\\
 		& \phantom{GRADING SCHEME:} \footnotesize {60\% of total module marks}   \\[0.8cm]
 		EXTERNAL EXAMINER: Prof. John King & \\[1cm]
 		\\[1cm]
 	\end{tabular}
 	\begin{center}
 		{\bf INSTRUCTIONS TO CANDIDATES}
 	\end{center}
 	
 	{\noindent \\ This paper comprises six questions. Question 1 is compulsory and is worth 40 Marks.\\  You must also attempt any four of the other five questions, each of which are worth 15 marks.
 		\\ Scientific calculators approved by the University of Limerick can be used. 
 		\\ Formula sheet and statistical tables are provided.
 	}
 	\normalsize
 	\newpage
 	\large
 	%\section*{Attempt ALL questions}
\subsection*{Question 1 - Short Questions (40 Marks)}

%========================================================%
% Short Questions
% 1 - Inverse
% 2 - Domain and Range
% 3 - Vertical Asymptotes
% 4 - Horizontal Asymptotes
% 5 - Limit ( Indeterminate Answer - Factorize)
%======================================================%
\begin{itemize}
	
	\item[(i)](4 Marks) Find $f^{-1}(x)$ the inverse of the function
	\[f(x) = e^{3x}. \]
		\smallskip
	\item[(ii)](4 Marks)	Solve the following limit:
	\[\lim_{x \to 2 } \frac{x^2- 7x +10}{x-2}.\]
	\smallskip
	\item[(iii)] (4 Marks) The domain and range of the function $f(x) = \cos(x)$ are $(-\infty, \infty)$ and $[-1,1]$ respectively. State the domain and the range of the following functions:
	\begin{multicols}{2}
		\begin{itemize}
		
		
	\item[(a)] $f(x) = \sin^2(x)  $,
		\item[(b)] $f(x) = \cos(3x)  $,
	\item[(c)] $f(x) = 3\sin(x)  $,
		\item[(d)] $f(x) = \sin(x) + 5$.
	\end{itemize}
	\end{multicols}
	\smallskip
	\item[(iv)] (4 Marks) Determine the vertical asymptote(s) of the following function:
	
	
	\[ y = f(x) = \frac{3x^2}{x^2-36}. \]	
	\smallskip
		\item[(v)] (4 Marks) Determine the horizontal asymptote(s) of the following function:
		
	\[ y = f(x) = \frac{3x^2}{x^2-36}. \]
		
		% answer = 3
		\smallskip
	\item[(vi)] (4 Marks) Find the x-intercepts and the y-intercept of the following function:
	
	\[ f(x) = x^2 + 3x - 18 . \]
	
	\smallskip


	%======================================================%
	% Short Questions
	% 6 - Indefinite Integral
	% 7 - Definite Integral
	% 8 - Partial Derivatives
	% 9 - Convergence
	% 10 - Sum of a Series
	\newpage
	%======================================================%
	\item[(vii)] (4 Marks) Evaluate the following indefinite integrals:
	\begin{multicols}{2}
	\begin{itemize}
		\item[(a)] $\displaystyle{ \int \sqrt{x} + e^{2x} + 1; dx, }  $
		\item[(b)] $\displaystyle{ \int \;\sin(2x)+ \sinh(x) + x \; dx}$
		
	\end{itemize}
	\end{multicols}

	\smallskip
%	\item (4 Marks) Evaluate the following indefinite integral:
%	\[ \int \cos(5x)+e^{2x}\; dx   \]
	%======================================================%
	\item[(viii)] (4 Marks)Evaluate the following definite integral:
	
	\[  \int^{4}_{1}  2x+5\; dx.  \]
	\smallskip

	
%	%======================================================%
%	\item[(viii)] Determine both of the first order partial derivatives with respect to $z$ of the following expression.
%	
%	\[ z = \frac{x^2y^3}{3} +  y \cos(x)\]
	
	
	\item[(ix)] (4 Marks) Find $\displaystyle{ \frac{ \partial^2 z }{ \partial x^2 }}$ and $\displaystyle{ \frac{ \partial^2 z }{ \partial y^2 }}$  for both of the following functions:
%	\[ z = 3x^2y^3 + y\cosh(x) - 9y^2 +10x^3.\]	
\smallskip
	\begin{multicols}{2}
		\begin{itemize}
			\item[(a)] $\displaystyle{  z = 6y^2 + 12x^3. }  $
				\item[(b)] $\displaystyle{  z = 2x^2y^3 + y\cosh(x).}$
				
			\end{itemize}
		\end{multicols}
		\medskip
	%========================================================%
	% READY
	\item[(x)] (4 Marks)  Compute the following summation:
	
	\[ \sum_{i=20}^{60} i . \]
\end{itemize}





	

	
%		
%
%		\item (4 Marks) Find the sum of the following geometric series: 
%		\[2 + 6 + 18 +  \ldots + 39366		\]
%
%		\item (4 Marks) Express the following repeating decimal number as a simple fraction. Show your workings.
%		
%		\[0.297297297....\]
%		\textit{(Simplify your answer such that both the numerator and the denominator are prime numbers.)}
%
%		\item (4 Marks) The three terms below are three successive terms in an aritmetic progression. Compute the value for $x$, and the common difference $d$
%		\[ \ldots,\; 3x-1,\; 2x+6,\;	5x-7,\; \ldots		  \]

\bigskip
\subsection*{Question 2 - Limits and Functions (15 Marks)}
\subsubsection*{Part A - Limits (5 Marks)}
\begin{enumerate}
	\item[(i)] (2 Marks)  Compute the limit of the following function.
		{
			\large
	\[\lim_{x \to 6 } \frac{x^2 + 2x-10}{x-4}.\]
}
	\bigskip
	
	\item[(iii)] (3 Marks) Compute the limit of the following function.
		{
			\large
	\[ \lim_{x \to \infty } \frac{3 + 2x^2 - 6x^3 }{4x^3 - 5x + 7}. \]	
}
\end{enumerate}
\medskip
\subsubsection*{Part B - Functions (6 Marks)}
\begin{enumerate}[(i)]
	\item (3 Marks) Determine if the function $f(x) = x^3\sin(x)$ is an even function, an odd function or neither.
	\bigskip
		\item (3 Marks) Determine if the function $f(x) = x^2+x$ is an even function, an odd function or neither.
		\bigskip
%	\item (2 Marks) Given the functions $g(x) = x^2+1$ and $f(x) = 1-3x$ determine expressions for $f \circ g(x)$ and $g \circ f(x)$.
\end{enumerate}
\medskip
%	Odd or Even functions %DONE IN CLASS LECTURE 6A%
%	\[f(x) = x^4\sin(x)\]
%	\[f(x) = x^3\cos(x)\]
\subsubsection*{Part C - Hyperbolic Functions (4 Marks)}


\begin{enumerate}[(i)]
	\item (4 Marks) Using their definition in terms of exponentials, prove the following hyperbolic identity   
	{
		\large
	\[\sinh^2 (x) = \frac{1}{2}\left[\cosh(2x)-1\right]\]
}
	%	\item Show that \[Cosh^2 x = Cosh2x + Sinh2x\]	
	%	\item Show that \[ cosh(x+y) = cosh(x)cosh(y) + sinh(x)sinh(y)\]
	
\end{enumerate}
%========================================================================================= %
%========================================================================================= %

\bigskip
\subsection*{Question 3 - Curve Sketching (15 Marks)}
Consider the function $f(x) = x^4 -  8x^2 + 9$

\begin{enumerate}[(i)]
	\item (2 Marks) Find the y intercept of the function $f(x)$.
	\item (3 Marks) Find the three turning points of the function $f(x)$ and classify them as local
	maxima or local minima.
	\item (3 Marks)  Find the two points of inflection of the function $f(x)$.
	
	\item (3 Marks)  Determine the behaviour of the function $f(x)$ as $x \rightarrow + \infty$ and as $x \rightarrow - \infty $.
	\item (4 Marks) Sketch the graph of the function $f(x)$ illustrating the features of the curve obtained
	in parts (i – iv). 
\end{enumerate}



\subsection*{Question 4 - Sequences and Series (15 Marks)}
% Sequences and Series

% 4 Marks - Geometric Series
% 3 Marks - Arithmetic Series
% 3 Marks - Telescoping Series
% 3 Marks - Ratio Test
% 2 Marks - Repeating Decimal 

\begin{enumerate}[(i)]
	\item (3 Marks)
	The second term $u_2$ of a geometric sequence is 4. The fifth term $u_5$ is -32. \\ \bigskip Answer the following questions. Both questions are worth 2 Marks each.
	\begin{itemize}
		\item[(a)] Find the common ratio $r$. 
		\item[(b)] Find the third and sixth term $u_3$ and $u_6$.
	\end{itemize}
	
	
	\bigskip
	
	\item (3 Marks)	Three consecutive terms of an arithmetic series are 
	\[4x+3,\;\;6x-8,\;\;8x-19.\]
	Find the values for $x$ and constant difference $d$.
	%answer is 460
	
		\bigskip
		
		\item (2 Marks) Express the following repeating decimal number as a simple fraction. Show your workings.
		
		\[0.513513513513....\]
		
	\bigskip
	%------------------------%
	%%updated for 2016
	\item (4 Marks) Suppose that the following term is the general term for a series. Use the Ratio Test to test this series for convergence
	
	\[u_n=\frac{2^n}{n!}\]
	
	\bigskip
	
	%------------------------%
	%%updated for 2016
	\item (3 Marks) Find the sum of the telescoping series  \[ \sum^{\infty}_{n=1} \frac{2
		}{(2n-1)(2n+1)}\]
	

	%	\item (2 Marks) Find the sum of the following telescoping series
	%	\[  \sum^{\infty}_{n=1}   \frac{6}{(3n+1)(3n+4)}  \]
	%	
	%	
	\newpage

\end{enumerate}
\subsection*{Question 5 - Integration (15 Marks)}


%Integration
%Substitution Technique
%Area Between Curves
%Definite Integrals

% - http://media.wix.com/ugd/b064dd_62ab7d7000c34b7eb4a8e2108269cf27.pdf

\begin{enumerate}[(i)]
	\item (3 Marks)  Evaluate the definite integral
	\[ \int^3_1 4x^2 + \sqrt[3]{x} + \frac{e^x}{3} \;dx. \]
	% \[ \int cos x(1 + sin x)^3 dx \]
	%---------------------------%
	\item (3 Marks) By finding a good substitution, evaluate the indefinite integral
	
	\[ \int \frac{18x + 18}{3x^2 + 6x + 7} \;dx.\]
	
	\item (3 Marks) By finding a good substitution, evaluate the indefinite integral
	% - \[ \int 3x^2\sqrt{4x^3 + 2}\; dx.\]
	% - \[ \int x(x^2 − 1)^5 dx\]
	\[ \int 2x \cosh(3x^2) dx \]
	% - \[ \int \frac{12x + 14}{3x^2 + 7x + 3} dx\]
	%---------------------------%
	\item (3 Marks) 
	%Use integration by parts to evaluate the indefinite integral
	%\[ \int  3xe^{3x} \;dx. \]
   Use integration by parts to find \[ \int x \cos(x) dx \]
	%---------------------------%
	\item (3 Marks) By first performing a partial fraction expansion (that is, by writing the integrand
	as follows)
	\[  \frac{A}{x - 1} + \frac{B}{x + 2},
	\]evaluate the definite integral
	
	\[ \int \frac{x + 11}{x^2 - x - 2}  dx. \]
	
\end{enumerate}
\bigskip 

%========================================================================================= %
\newpage
\subsection*{Question 6 - Applications of Calculus (15 Marks)}

\subsubsection*{Part A - Applications of Integration}
\begin{enumerate}[(i)]
	\item (5 Marks) A current $i(t) = 4 + 4\cos(5t)$ passes through a capacitor at time $t$.
	The capacitor is uncharged at time $ t = 0 $. Find the charge $q(t)$ at all times $t$.
	%\item (5 Marks) A moving object has acceleration $a(t) = 3 + 5 sin t$ at time t. It starts from rest at time $t = 0$. Find its velocity at all time t.	Also, find its velocity at time $t = 4$.
	\smallskip
	\item (5 Marks) Find the area enclosed by the curve $y=1+2x-x^2$ and the line $y=x-1$.
	Clearly state the points of intersection in your answer.
	\smallskip
	%\item (5 Marks) Find the area enclosed by the curve $y = 3x^2 − 4x + 1$ and the x axis.

\end{enumerate}



%\item Use Simpsons with 5 equal subintervals to find an approximation
%for
\bigskip
\subsubsection*{Part B - Partial Derivatives}
\begin{enumerate}[(i)]
	
	\item (5 Marks)
	Show that the function $z=e^{-4y}\cos(x)$ satisfies the partial differential equation
	\[ 16\frac{\partial^2 z}{\partial x^2} + \frac{\partial^2 z}{\partial y^2} = 0. \]
\end{enumerate}
%
%\subsubsection*{MacLaurin Series}
%\begin{enumerate}[(i)]
%	\item	Find the Maclaurin series of $e^x$ up to an including the term containing $x^4$
%	
%	%\item Evaluate all the first partial derivatives, i.e. $\frac{\partial z}{\partial x}$ and $\frac{\partial z}{\partial y} $, 
%	%of the following function:
%	%
%	%\[ z = 2xy^3 + cos(2x) - sin(y) \]
%
%
%
%	\item
%\end{enumerate}

\newpage
	\section*{Formula Sheet}
	
\subsection*{Logarithms}
If $a^b = c$ \, then \, $\mbox{log}_a c = b$.



\subsection*{Change of Base Formula}

\[ \log_A(B) = \frac{ \log_e(B) }{ \log_e(A) }  \]
	\subsection*{Sum and Difference of Two Cubes}
	\[ a^3 + b^3 = (a-b)(a^2 - ab + b^2)\]
	\[ a^3 - b^3 = (a-b)(a^2 + ab + b^2)\]
	
	%======================================== %
	
	\subsection*{Sequences and Series}
Arithmetic Series Summation:
\begin{multicols}{2}	
	\[ \sum_{i=1}^{n} i = \frac{n(n+1)}{2}\]
	

	\[ S_n = \frac{n}{2} \left(2a + (n-1) d \right)\]
\end{multicols}	
	Geometric Series Summation:
	\begin{multicols}{2}
	\[ S_n = a\left(\frac{1-r^n}{1-r}\right)\]
	
	\[ S_\infty = \frac{a}{1-r}\]
	\end{multicols}
	
	
	\subsection*{Ratio Test}
	
	For a series with general term $u_n$, if
	
	\[ \lim_{n \to \infty } \left| \frac{u_{n+1}}{u_n} \right| = r\]
	then the series converges (absolutely) if $r<1$
	
	
	%==========================================================================================%
	
	
	\subsection*{Curve Sketching}
	\begin{description}
		\item[Horizontal Asymptote:] The horizontal asymptote is computed as
		\[ \lim_{x \to \infty } f(x) \]
	\end{description}
	%==========================================================================================%
	
	\subsection*{Maclaurin Series}
	\[f(x) = f(0) + f^{\prime}(0) + \frac{f^{\prime \prime}(0)}{2!} + \frac{f^{\prime\prime \prime}(0)}{3!} + \ldots \]
	%==========================================================================================%
	
	\subsection*{Hyperbolic Functions }
	
	\begin{multicols}{2}
	\[ \cosh(x)  =  \frac{e^{x} + e^{-x}}{2} \]
	
	\[ \sinh(x)  = \frac{e^{x} - e^{-x}}{2} \]
	\end{multicols}

	
	%==========================================================================================%
	

	\subsection*{Integration}
	
	Integration by parts: 
	
	\[ \int u dv = uv - \int v du \]  
	
\noindent Further formulae and special cases on pages 25 \& 26 of the log tables provided.
	
	\subsection*{Dynamics}
	Where $s(t)$ denotes displacement at time $t$, $v(t)$ denotes the velocity at time $t$ and $a(t)$
	denotes the acceleration at time $t$, 
	\begin{multicols}{2}
	\[  \frac{ds(t)}{dt}  = v(t),\]
	\[  \frac{dv(t)}{dt}  = a(t).\]
	\end{multicols}
	\subsection*{Electrical Circuits}
	Where $q(t)$ denotes the charge at time $t$ and $i(t)$ denotes the current at time $t$,
	\[  \frac{dq(t)}{dt}  = i(t).\]
	
	

\end{document}

	\subsection*{Rules of Differentiation}
	%% - http://media.wix.com/ugd/b064dd_62ab7d7000c34b7eb4a8e2108269cf27.pdf
	\begin{description}
		\item[Product Rule:]  with $y = uv$
		
		
		\[ \frac{dy}{dx} = u \frac{dv}{dx} +  v \frac{du}{dx} \]
		
		\item[Quotient Rule:] \[ y = \frac{u}{v}\]
		\[ \frac{dy}{dx}  = \frac{v \frac{du}{dx} - u \frac{dv}{dx} }{v^2} \]
		
		
		
		\item[Chain Rule:]
		
		% y = f(u) and u = u(x), that is y = f(u(x)) ⇒ dy
		% dx
		\[ \frac{dy}{dx} = \frac{dy}{du} \times \frac{du}{dx}  \]
	\end{description}
	%==========================================================================================%
